% Modul Belajar 5, dihasilkan oleh tools/build.py dari tools/konten/p05.py
\documentclass[a4paper]{article}
\usepackage{modulITERA}
\begin{document}
\Sampul{5}{Perulangan Bagian 1}{for, while, dan do-while dengan pola pencacah, akumulator, dan sentinel}{CPMK0607 (menerapkan) dan CPMK0608 (menjelaskan) pada konsep dasar algoritma dan sistem komputer}{sekitar 3 jam belajar mandiri, ditambah 150 menit di kelas}{\texttt{02 Hands-on/P5 - Perulangan 1 - Hands-on/}}
\Bagian{Modul 5}{Peta belajar}
Ikuti urutan ini dari atas ke bawah. Setiap bagian memakai apa yang sudah dibahas di bagian sebelumnya.
\par\vspace{4pt}\noindent{\fontsize{9.5pt}{13pt}\selectfont
\begin{tabular}{@{}>{\raggedright\arraybackslash}p{0.140\textwidth}>{\raggedright\arraybackslash}p{0.840\textwidth}@{}}
\kepala{Urutan} & \kepala{Bagian}\\
1 & Tentang modul ini\\\garis
2 & Masalah pembuka: Rata-rata nilai satu kelas\\\garis
3 & Langkah 1: Perulangan for\\\garis
4 & Langkah 2: Pencacah dan akumulator\\\garis
5 & Langkah 3: Perulangan while\\\garis
6 & Langkah 4: Nilai sentinel\\\garis
7 & Langkah 5: Perulangan do-while\\\garis
8 & Langkah 6: Memilih jenis perulangan\\\garis
9 & Pesan error yang sering muncul\\\garis
10 & Latihan mandiri\\\garis
11 & Rangkuman\\\garis
12 & Cek pemahaman\\\garis
13 & Exit Ticket dan tugas\\\garis
14 & Bacaan lanjut dan persiapan minggu depan\\
\garisdasar
\end{tabular}}\par\vspace{6pt}
\Bagian{Pembuka}{Tentang modul ini}
Modul ini mendampingi Pertemuan 5. Program kita sudah bisa memilih. Minggu ini program belajar mengulang: menjalankan blok yang sama berkali-kali tanpa menulis ulang kodenya, sambil mengumpulkan hasil di setiap putaran.

Isinya sama dengan slide di kelas, tetapi ditulis lebih pelan dan lebih lengkap, supaya kalian bisa mengulang sendiri di rumah tanpa harus menebak maksud slide.

\Sub{Setelah menyelesaikan modul ini, kalian bisa}
\par\vspace{4pt}\noindent{\fontsize{9.5pt}{13pt}\selectfont
\begin{tabular}{@{}>{\raggedright\arraybackslash}p{0.070\textwidth}>{\raggedright\arraybackslash}p{0.680\textwidth}>{\raggedright\arraybackslash}p{0.230\textwidth}@{}}
\kepala{No} & \kepala{Kemampuan} & \kepala{Dibahas di}\\
1 & Menulis program yang mengulang proses dengan \texttt{for}, \texttt{while}, dan \texttt{do-while}, memakai pola pencacah, akumulator, dan nilai sentinel (Sub-CPMK 5.1, CPMK0607) & Langkah 1 sampai 6\\\garis
2 & Menelusuri perulangan dengan trace table, menentukan berapa kali badan perulangan dijalankan, dan menjelaskan penyebab off-by-one atau perulangan tak berhenti (Sub-CPMK 5.2, CPMK0608) & kotak Tebak dulu, Latihan 3\\
\garisdasar
\end{tabular}}\par\vspace{6pt}
Karena setiap instrumen penilaian dibagi rata antara Bagian A (menulis kode) dan Bagian B (menelusuri dan menjelaskan), yang dinilai bukan hanya program kalian jalan, tetapi juga kalian bisa menjelaskan mengapa program itu jalan.

\Sub{Cara memakai modul ini}
Setiap langkah punya pola yang sama: penjelasan konsep, kadang contoh dari kehidupan sehari-hari, lalu kode yang bisa langsung kalian ketik. Sesudah kode ada kotak \textbf{Tebak dulu}. Tulis tebakan kalian di kertas sebelum membaca \textbf{Jawaban} di bawahnya. Kebiasaan menebak lalu membuktikan inilah yang membuat konsep menempel, bukan sekadar membaca.

\begin{Penting}[Ketik, jangan salin]
Ketik ulang setiap contoh di GDB Online atau editor kalian, jangan disalin-tempel. Saat mengetik, kalian akan membuat salah ketik, membaca pesan error, dan memperbaikinya. Itu bagian dari belajar.
\end{Penting}
\Sub{Yang perlu disiapkan}
Peramban dengan akses ke onlinegdb.com, atau g++ di komputer kalian sendiri. Kertas dan alat tulis untuk menebak keluaran dan membuat trace table. Folder hands-on pertemuan ini dari repo mata kuliah.

\Bagian{Pembuka}{Masalah pembuka: Rata-rata nilai satu kelas}
Asisten praktikum harus menghitung rata-rata nilai Exit Ticket. Jumlah mahasiswa berbeda di setiap kelas, kadang 30, kadang 45.

Menulis \texttt{cin} sebanyak 45 kali jelas tidak masuk akal. Lagi pula banyaknya data baru diketahui saat program berjalan.

\begin{MKeluaran}[title={Tampilan yang diinginkan}]
Masukkan nilai (-1 untuk selesai):
70 85 60 90 -1
Banyak    : 4
Tertinggi : 90
Rata-rata : 76.25
\end{MKeluaran}
\begin{Tebak}[Coba pikirkan]
Langkah apa yang terus diulang untuk setiap nilai? Bagaimana program tahu kapan berhenti?
\end{Tebak}
\begin{Jawaban}[Cara memecahnya]
Temukan yang berulang: Baca satu nilai, tambahkan ke total, tambah pencacah. Tentukan kapan berhenti: Berhenti saat bertemu -1, atau setelah n kali. Terjemahkan: \texttt{while} untuk syarat berhenti, \texttt{for} untuk banyak putaran yang sudah diketahui.
\end{Jawaban}
\Bagian{Langkah 1}{Perulangan for}
\texttt{for (awal; syarat; perubahan)} cocok bila banyak putaran sudah diketahui. Urutannya: jalankan bagian awal sekali, periksa syarat, jalankan badan, jalankan perubahan, periksa syarat lagi, dan seterusnya sampai syaratnya salah.

Variabel \texttt{i} yang dideklarasikan di dalam \texttt{for} hanya hidup di dalam perulangan itu.

\begin{MKode}{for.cpp}
for (int i = 1; i <= 5; i++) {
    cout << i << " ";
}
cout << endl;
for (int i = 10; i > 0; i -= 3) {
    cout << i << " ";
}
cout << endl;
\end{MKode}
\begin{Tebak}[]
Sebelum menjalankan program, tulis keluarannya di kertas, karakter demi karakter.
\end{Tebak}
\begin{Jawaban}[]
\begin{MKeluaran}[]
1 2 3 4 5 
10 7 4 1 
\end{MKeluaran}
Tiga bagian \texttt{for}: nilai awal; syarat lanjut; perubahan. Syarat diperiksa sebelum setiap putaran.
\end{Jawaban}
\Bagian{Langkah 2}{Pencacah dan akumulator}
Dua pola paling sering dipakai. Akumulator mengumpulkan nilai (\texttt{jumlah += i}), pencacah menghitung kejadian (\texttt{genap++}). Keduanya harus diberi nilai awal sebelum perulangan. Bila diberi nilai awal di dalam perulangan, nilainya ter-reset di setiap putaran.

\begin{MKode}{akumulator.cpp}
int jumlah = 0, genap = 0;
for (int i = 1; i <= 10; i++) {
    jumlah += i;
    if (i % 2 == 0) {
        genap++;
    }
}
cout << "Jumlah: " << jumlah << endl;
cout << "Genap : " << genap << endl;
\end{MKode}
\begin{Tebak}[]
Sebelum menjalankan program, tulis keluarannya di kertas, karakter demi karakter.
\end{Tebak}
\begin{Jawaban}[]
\begin{MKeluaran}[]
Jumlah: 55
Genap : 5
\end{MKeluaran}
Akumulator dan pencacah selalu diberi nilai awal sebelum perulangan, bukan di dalamnya.
\end{Jawaban}
\Bagian{Langkah 3}{Perulangan while}
\texttt{while (syarat)} cocok bila banyak putaran belum diketahui, hanya syarat berhentinya. Syarat diperiksa sebelum setiap putaran, sehingga badan bisa saja tidak dijalankan sama sekali.

Pastikan ada pernyataan di dalam badan yang membuat syaratnya suatu saat salah. Bila lupa, perulangan tidak berhenti. Di GDB Online, tekan Stop untuk menghentikannya.

\begin{MKode}{while.cpp}
int saldo = 100000, bulan = 0;
while (saldo < 500000) {
    saldo += 150000;
    bulan++;
}
cout << bulan << " bulan, saldo " << saldo << endl;
\end{MKode}
\begin{Tebak}[]
Sebelum menjalankan program, tulis keluarannya di kertas, karakter demi karakter.
\end{Tebak}
\begin{Jawaban}[]
\begin{MKeluaran}[]
3 bulan, saldo 550000
\end{MKeluaran}
Sesuatu di dalam badan harus mengubah syarat. Bila tidak, perulangan tidak pernah berhenti.
\end{Jawaban}
\Bagian{Langkah 4}{Nilai sentinel}
Sentinel adalah nilai khusus yang menandai akhir data, misalnya -1 untuk nilai ujian. Polanya: baca satu nilai sebelum perulangan, periksa di syarat \texttt{while}, proses, lalu baca nilai berikutnya di akhir badan. Urutan ini menjamin sentinel tidak ikut diproses.

\begin{MKode}{sentinel.cpp}
int nilai, banyak = 0, total = 0;
cin >> nilai;
while (nilai != -1) {
    banyak++;
    total += nilai;
    cin >> nilai;
}
cout << banyak << " data, total " << total << endl;
\end{MKode}
\begin{MKeluaran}[title={Masukan}]
70 85 60 90 -1
\end{MKeluaran}
\begin{Tebak}[]
Sebelum menjalankan program, tulis keluarannya di kertas, karakter demi karakter.
\end{Tebak}
\begin{Jawaban}[]
\begin{MKeluaran}[]
4 data, total 305
\end{MKeluaran}
Baca sekali sebelum perulangan, lalu baca lagi di akhir badan. Dengan begitu -1 tidak ikut dihitung.
\end{Jawaban}
\Bagian{Langkah 5}{Perulangan do-while}
\texttt{do \{ ... \} while (syarat);} menjalankan badan dulu, baru memeriksa syarat. Badannya pasti jalan minimal sekali, sehingga cocok untuk menu dan validasi masukan.

\begin{MKode}{dowhile.cpp}
int pin, coba = 0;
do {
    cin >> pin;
    coba++;
} while (pin != 1234 && coba < 3);
cout << (pin == 1234 ? "Masuk" : "Terkunci");
cout << " setelah " << coba << " kali" << endl;
\end{MKode}
\begin{MKeluaran}[title={Masukan}]
1111 1234
\end{MKeluaran}
\begin{Tebak}[]
Sebelum menjalankan program, tulis keluarannya di kertas, karakter demi karakter.
\end{Tebak}
\begin{Jawaban}[]
\begin{MKeluaran}[]
Masuk setelah 2 kali
\end{MKeluaran}
Syarat diperiksa sesudah badan. Jangan lupa titik koma setelah \texttt{while (...)}.
\end{Jawaban}
\Bagian{Langkah 6}{Memilih jenis perulangan}
Secara teknis ketiga perulangan bisa saling menggantikan. Pilihan yang tepat membuat maksud program langsung terbaca.

\par\vspace{4pt}\noindent{\fontsize{9.5pt}{13pt}\selectfont
\begin{tabular}{@{}>{\raggedright\arraybackslash}p{0.414\textwidth}>{\raggedright\arraybackslash}p{0.174\textwidth}>{\raggedright\arraybackslash}p{0.392\textwidth}@{}}
\kepala{Situasi} & \kepala{Pilihan} & \kepala{Contoh}\\
Banyak putaran diketahui & \texttt{for} & tampilkan 1 sampai n\\\garis
Berhenti karena syarat, bisa nol putaran & \texttt{while} & baca sampai sentinel\\\garis
Badan harus jalan minimal sekali & \texttt{do-while} & menu, validasi PIN\\
\garisdasar
\end{tabular}}\par\vspace{6pt}
\begin{Penting}[]
Ketiganya bisa saling menggantikan. Pilih yang paling jelas menyatakan maksud kalian.
\end{Penting}
\Bagian{Waspada}{Pesan error yang sering muncul}
Semua pesan di tabel ini dihasilkan g++ dengan opsi \texttt{-Wall}, sama seperti yang dipakai \texttt{cek.sh}.
\par\vspace{4pt}\noindent{\fontsize{9.5pt}{13pt}\selectfont
\begin{tabular}{@{}>{\raggedright\arraybackslash}p{0.240\textwidth}>{\raggedright\arraybackslash}p{0.270\textwidth}>{\raggedright\arraybackslash}p{0.470\textwidth}@{}}
\kepala{Kesalahan} & \kepala{Contoh} & \kepala{Pesan pertama dari g++}\\
Koma di header for & \texttt{for (int i = 0, i < 5, i++)} & \texttt{error: expected ';' before '<' token}\\\garis
Variabel di luar jangkauan & \texttt{cout << i; sesudah for} & \texttt{error: 'i' was not declared in this scope}\\\garis
Lupa ; sesudah do-while & \texttt{\} while (x < 5)} & \texttt{error: expected ';' before 'cout'}\\\garis
Perbandingan bertipe beda & \texttt{i < s.length() dengan int} & \texttt{warning: comparison of integer expressions of different signedness: 'int' and ...}\\\garis
Badan for kosong & \texttt{for (...);} & \texttt{warning: this 'for' clause does not guard...}\\
\garisdasar
\end{tabular}}\par\vspace{6pt}
Off-by-one (kelebihan atau kekurangan satu putaran) dan perulangan tak berhenti tidak dilaporkan compiler. Keduanya hanya ketahuan lewat trace table atau pengujian, karena itu latihan tantangan minggu ini adalah menelusuri perulangan.

\Bagian{Latihan}{Latihan mandiri}
Latihan ini sama dengan hands-on di kelas. Semua file ada di \texttt{02 Hands-on/P5 - Perulangan 1 - Hands-on/latihan/}. Kerjakan berurutan, lalu cocokkan dengan \texttt{expected/} atau jalankan \texttt{bash cek.sh}.
\par\vspace{4pt}\noindent{\fontsize{9.5pt}{13pt}\selectfont
\begin{tabular}{@{}>{\raggedright\arraybackslash}p{0.070\textwidth}>{\raggedright\arraybackslash}p{0.360\textwidth}>{\raggedright\arraybackslash}p{0.150\textwidth}>{\raggedright\arraybackslash}p{0.400\textwidth}@{}}
\kepala{No} & \kepala{File} & \kepala{Tingkat} & \kepala{Yang dilatih}\\
1 & \texttt{handson1-deret.cpp} & Dasar & \texttt{for}, akumulator\\\garis
2 & \texttt{handson2-tabungan.cpp} & Menengah & \texttt{while} dengan syarat berhenti\\\garis
3 & \texttt{handson3-rerata.cpp} & Tantangan & trace table, perbaiki tiga kesalahan sentinel\\\garis
+ & \texttt{bonus-faktorial.cpp} & Pengayaan & \texttt{do-while} untuk menu, \texttt{long long}\\
\garisdasar
\end{tabular}}\par\vspace{6pt}
\Sub{Latihan 1: handson1-deret.cpp}
Tampilkan 1 sampai n dan jumlahnya. Masukan uji: \texttt{6}.

\Sub{Latihan 2: handson2-tabungan.cpp}
Simulasikan saldo tabungan per bulan sampai mencapai target.

\Sub{Latihan 3: handson3-rerata.cpp}
Program rata-rata dengan sentinel -1 memberi hasil salah. Buat trace table, perbaiki, dan jelaskan.

\Sub{Bonus: bonus-faktorial.cpp}
Menu berulang untuk faktorial dan jumlah deret sampai pengguna memilih 0.

\Sub{Latihan menelusuri}
\begin{MKode}{tebak.cpp}
int x = 0;
for (int i = 0; i < 4; i++) {
    x += i * 2;
}
int k = 1;
while (k < 20) {
    k *= 3;
}
cout << x << " " << k << endl;
\end{MKode}
\begin{Tebak}[]
Tulis keluarannya di kertas tanpa menjalankan program. Buat trace table bila perlu.
\end{Tebak}
\begin{Jawaban}[]
\begin{MKeluaran}[]
12 27
\end{MKeluaran}
\texttt{x} = 0 + 2 + 4 + 6 = 12. \texttt{k} menjadi 3, 9, 27; berhenti saat 27 tidak lagi kurang dari 20.
\end{Jawaban}
\Bagian{Penutup}{Rangkuman}
\par\vspace{4pt}\noindent{\fontsize{9.5pt}{13pt}\selectfont
\begin{tabular}{@{}>{\raggedright\arraybackslash}p{0.320\textwidth}>{\raggedright\arraybackslash}p{0.660\textwidth}@{}}
\kepala{Konsep} & \kepala{Yang perlu diingat}\\
Perulangan for & Tiga bagian \texttt{for}: nilai awal; syarat lanjut; perubahan.\\\garis
Pencacah dan akumulator & Akumulator dan pencacah selalu diberi nilai awal sebelum perulangan, bukan di dalamnya.\\\garis
Perulangan while & Sesuatu di dalam badan harus mengubah syarat.\\\garis
Nilai sentinel & Baca sekali sebelum perulangan, lalu baca lagi di akhir badan.\\\garis
Perulangan do-while & Syarat diperiksa sesudah badan.\\\garis
Memilih jenis perulangan & Ketiganya bisa saling menggantikan.\\
\garisdasar
\end{tabular}}\par\vspace{6pt}
\Bagian{Penutup}{Cek pemahaman}
Jawab tanpa menjalankan program, lalu periksa dengan menjalankannya.
\begin{enumerate}
\item Berapa kali badan \texttt{for (int i = 2; i < 10; i += 2)} dijalankan?
\item Mengapa \texttt{int total = 0;} harus di luar perulangan?
\item Apa beda \texttt{while} dan \texttt{do-while} bila syaratnya sejak awal salah?
\item Apa keluaran \texttt{int k = 1; while (k < 50) k *= 2; cout << k;}?
\item Pada pola sentinel, mengapa \texttt{cin} ditulis di akhir badan?
\end{enumerate}
\begin{Jawaban}[]
\begin{enumerate}[leftmargin=16pt]
\item 4 kali (i = 2, 4, 6, 8).
\item Bila di dalam, total di-reset di setiap putaran.
\item \texttt{while} tidak menjalankan badan; \texttt{do-while} menjalankannya sekali.
\item 64.
\item Agar sentinel diperiksa di syarat sebelum sempat diproses.
\end{enumerate}
\end{Jawaban}
\Bagian{Penilaian}{Exit Ticket 5}
Dikerjakan di 25 menit terakhir kelas, sendiri, tanpa AI, internet, atau diskusi. Exit Ticket 5 adalah satu dari 14 Exit Ticket; rata-ratanya menjadi komponen berbobot 25\%.
\par\vspace{4pt}\noindent{\fontsize{9.5pt}{13pt}\selectfont
\begin{tabular}{@{}>{\raggedright\arraybackslash}p{0.080\textwidth}>{\raggedright\arraybackslash}p{0.600\textwidth}>{\raggedright\arraybackslash}p{0.100\textwidth}>{\raggedright\arraybackslash}p{0.200\textwidth}@{}}
\kepala{Soal} & \kepala{Isi} & \kepala{Poin} & \kepala{CPMK}\\
A1 & Tulis program yang menampilkan kelipatan k dari 1 sampai n beserta banyaknya & 25 & CPMK0607\\\garis
A2 & Tulis program yang membaca bilangan sampai 0 lalu menampilkan jumlah bilangan positif & 25 & CPMK0607\\\garis
B1 & Buat trace table sebuah perulangan dan tentukan keluarannya & 20 & CPMK0608\\\garis
B2 & Jelaskan mengapa sebuah perulangan tidak berhenti dan cara memperbaikinya & 20 & CPMK0608\\\garis
B3 & Tentukan berapa kali badan perulangan dijalankan dan jelaskan alasannya & 10 & CPMK0608\\
\garisdasar
\end{tabular}}\par\vspace{6pt}
\Bagian{Penilaian}{Tugas 5: Program Statistik Nilai}
Kumpulkan sebelum Pertemuan 6 lewat kanal pengumpulan yang diumumkan dosen.

\Sub{Bagian A: program (50 poin)}
Buat program statistik nilai dengan ketentuan berikut.
\begin{enumerate}
\item Baca jumlah siswa n, ulangi sampai n > 0
\item Baca nilai n siswa; setiap nilai harus 0 sampai 100, ulangi bila tidak valid
\item Tampilkan nilai tertinggi, nilai terendah, dan rata-rata kelas
\item Tampilkan jumlah siswa LULUS (nilai 75 ke atas) dan TIDAK LULUS
\item Menu \texttt{do-while} untuk mengulang atau keluar
\end{enumerate}
\begin{Contoh}[Bonus, tidak wajib]
Tampilkan persentase kelulusan dengan dua desimal.
\end{Contoh}
\Sub{Bagian B: penjelasan (50 poin)}
Komentar maksimal 150 kata dan trace table untuk n = 3 dengan nilai 80, 70, 90. Jelaskan alasan memilih \texttt{for}, \texttt{while}, atau \texttt{do-while} di setiap bagian, dan satu kesalahan off-by-one atau perulangan tak berhenti yang kalian temui.

\Sub{Rubrik Tugas 5}
\par\vspace{4pt}\noindent{\fontsize{8.5pt}{11.5pt}\selectfont
\begin{tabular}{@{}>{\raggedright\arraybackslash}p{0.190\textwidth}>{\raggedright\arraybackslash}p{0.070\textwidth}>{\raggedright\arraybackslash}p{0.200\textwidth}>{\raggedright\arraybackslash}p{0.180\textwidth}>{\raggedright\arraybackslash}p{0.170\textwidth}>{\raggedright\arraybackslash}p{0.170\textwidth}@{}}
\kepala{Kriteria} & \kepala{Poin} & \kepala{Sangat baik 85–100} & \kepala{Baik 70–84} & \kepala{Cukup 55–69} & \kepala{Kurang <55}\\
A1 Program berjalan & 20 & tanpa error dan peringatan & ada peringatan & perlu satu perbaikan & tidak terkompilasi\\\garis
A2 Validasi dan perulangan & 15 & validasi n dan nilai, menu berulang & satu validasi kurang & dua kurang & tanpa validasi\\\garis
A3 Ketepatan statistik & 15 & semua tepat termasuk n = 1 & satu salah & dua salah & sebagian besar salah\\\garis
B1 Trace table dan alasan & 30 & trace lengkap dan alasan tepat & satu baris trace salah & trace tidak lengkap & tidak ada atau disalin\\\garis
B2 Cerita error dan pengujian & 20 & pesan, sebab, perbaikan, uji & pesan tidak dikutip & hanya perbaikan & tidak ada\\
\garisdasar
\end{tabular}}\par\vspace{6pt}
Nilai kriteria = poin × skor level. Contoh: A1 di level Baik dengan skor 80 memberi 20 × 0,80 = 16 poin.

\Bagian{Penutup}{Bacaan lanjut dan persiapan minggu depan}
\begin{enumerate}
\item Deitel, P. dan Deitel, H. C++ How to Program, edisi ke-10, Bab 4 dan 5.
\item learncpp.com, Bab 8 (loops).
\item Gaddis, T. Starting Out with C++, Bab 5.
\item Slide \texttt{01 Slide/P5 Perulangan Bagian 1.pdf} dan hands-on di \texttt{02 Hands-on/P5 - Perulangan 1 - Hands-on/}.
\end{enumerate}
\begin{Penting}[Minggu depan]
Pertemuan 6 membahas perulangan bersarang, \texttt{break}, \texttt{continue}, dan pola bintang, sejalan dengan Algoritma Pertemuan 6.
\end{Penting}
\end{document}
